\subsection{Poincaré groupoid}

DEFINITION 1. --- Let X be a topological space, let \(c_{0}\) and \(c_{1}\) be paths in X and let \(\sigma\colon\mathbf{I}\times\mathbf{I}\to\mathbf{X}\) be a homotopy connecting \(c_{0}\) to \(c_{1}\) . One says that \(\sigma\) is a strict homotopy if the maps \(s\mapsto\sigma(0,s)\) and \(s\mapsto\sigma(1,s)\) are constant.

We say that paths \(c_{0}\) and \(c_{1}\) in X are strictly homotopic if there exists a strict homotopy connecting \(c_{0}\) to \(c_{1}\) .

Two strictly homotopic paths have the same origin and the same endpoint.

Example. --- Let c be a path in X and let \(\varphi\colon I\to I\) be a continuous map such that \(\varphi(0)=0\) and \(\varphi(1)=1\) . The paths c and \(c\circ\varphi\) are strictly homotopic. Indeed, the map \(\sigma\colon I\times I\to X\) defined by \(\sigma(t,s)=c((1-s)t+s\varphi(t))\) is a strict homotopy connecting c to \(c\circ\varphi\) .

Let X be a topological space. Recall (cf.~III, p.~257) that \(\Lambda(X)\) denotes the topological space \(\mathcal{C}_c(\mathbf{I};\mathrm{X})\) of paths in X and that for \(x\) , \(y \in X\) , \(\Lambda_{x,y}(X)\) is the subspace of \(\Lambda(X)\) formed by the paths with origin \(x\) and endpoint \(y\) . The family of sets \(\Lambda_{x,y}(X)\) , for \(x\) , \(y \in X\) , is a partition of the space of paths of X. By the canonical bijection (III, p.~257, remark 2) of \(\mathcal{C}(\mathbf{I} \times \mathbf{I};\mathrm{X})\) onto \(\mathcal{C}(\mathbf{I};\Lambda(\mathrm{X}))\) , strict homotopies correspond to paths \(c: \mathbf{I} \to \Lambda(X)\) whose image is contained in a subspace of the form \(\Lambda_{x,y}(\mathrm{X})\) . The relation “the paths \(c_0\) and \(c_1\) are strictly homotopic” is therefore an equivalence relation in \(\Lambda(\mathrm{X})\) (III, p.~259, prop. 5) and the equivalence classes for this relation are the arcwise connected components of the subspaces \(\Lambda_{x,y}(\mathrm{X})\) of the space of paths of X. We denote by \(\varpi_{x,y}(\mathrm{X})\) the set \(\pi_0(\Lambda_{x,y}(\mathrm{X}))\) and call a class of paths joining \(x\) to \(y\) any element of \(\varpi_{x,y}(\mathrm{X})\) .

DEFINITION 2. --- We call the loop space of X, and denote by \(\Omega(X)\) , the subspace of \(\Lambda(X)\) consisting of the loops (III, p.~256, def. 1) in X.

We denote \(\Omega_x(\mathrm{X})\) the set \(\Lambda_{x,x}(\mathrm{X})\) . The elements of \(\Omega_x(\mathrm{X})\) are called the loops in X at \(x\) and the elements of \(\varpi_{x,x}(\mathrm{X})\) are called loop classes in X at \(x\) . The map \(e_x: \mathbf{I} \to \mathbf{X}\) is constant with image \(x\) and is a loop, called the constant loop at \(x\) ; its strict homotopy class is denoted by \(\varepsilon_x\) . The map \(x \mapsto e_x\) from X to \(\Lambda(\mathrm{X})\) is continuous (III, p.~257, prop. 1).

Let X be a topological space and let x, y, z be points of X. By passing to arcwise connected components, we deduce from the continuous map \(c \mapsto \overline{c}\) from \(\Lambda_{x,y}(X)\) to \(\Lambda_{y,x}(X)\) (III, p.~258, corollary) a map from \(\varpi_{x,y}(X)\) to \(\varpi_{y,x}(X)\) , denoted \(\gamma \mapsto \overline{\gamma}\) . If \(\gamma \in \varpi_{x,y}(X)\) , \(\overline{\gamma}\) is called the inverse of the path class \(\gamma\) .

Similarly, if we identify the sets \(\pi_0(\Lambda_{x,y}(\mathrm{X})) \times \pi_0(\Lambda_{y,z}(\mathrm{X}))\) and \(\pi_0(\Lambda_{x,y}(\mathrm{X}) \times \Lambda_{y,z}(\mathrm{X}))\) (III, p.~260, prop. 6), we deduce from the continuous map \((c,d) \mapsto c * d\) from \(\Lambda_{x,y}(\mathrm{X}) \times \Lambda_{y,z}(\mathrm{X})\) to \(\Lambda_{x,z}(\mathrm{X})\) (III, p.~258, corollary), by passing to arcwise connected components, a map \(C_{x,y,z} \colon \varpi_{x,y}(\mathrm{X}) \times \varpi_{y,z}(\mathrm{X}) \to \varpi_{x,z}(\mathrm{X})\) . For \(\gamma \in \varpi_{x,y}(\mathrm{X})\) and \(\delta \in \varpi_{y,z}(\mathrm{X})\) , write \(\gamma \delta\) for the strict homotopy class \(C_{x,y,z}(\gamma,\delta)\) . It is called the composite of juxtaposable path classes \(\gamma\) and \(\delta\) .

We have \(\overline{\overline{\gamma}}=\gamma\) and \(\overline{\gamma\delta}=\overline{\delta}\overline{\gamma}\) .

PROPOSITION 1. --- Let X be a topological space, let x, y, z, u be points of X, and let \(\gamma_{1} \in \varpi_{x,y}(X)\) , \(\gamma_{2} \in \varpi_{y,z}(X)\) , \(\gamma_{3} \in \varpi_{z,u}(X)\) be path classes. We have

\begin{enumerate}
\def\labelenumi{(\arabic{enumi})}
\tightlist
\item
\end{enumerate}

\[
\varepsilon_ {x} \gamma_ {1} = \gamma_ {1} \varepsilon_ {y} = \gamma_ {1},\tag{2}
\]

\[
\gamma_ {1} \overline {{\gamma}} _ {1} = \varepsilon_ {x}, \qquad \overline {{\gamma}} _ {1} \gamma_ {1} = \varepsilon_ {y},\tag{3}
\]

\[
(\gamma_ {1} \gamma_ {2}) \gamma_ {3} = \gamma_ {1} (\gamma_ {2} \gamma_ {3}).
\]

Let \(c_{1} \in \Lambda_{x,y}(\mathrm{X})\) , \(c_{2} \in \Lambda_{y,z}(\mathrm{X})\) , \(c_{3} \in \Lambda_{z,u}(\mathrm{X})\) be representatives of \(\gamma_{1}\) , \(\gamma_{2}\) and \(\gamma_{3}\) respectively. Let \(\varphi: \mathbf{I} \to \mathbf{I}\) be the function defined by

\[
\varphi (t) = \left\{ \begin{array}{l l} t / 2 & \text { for } 0 \leqslant t \leqslant 1 / 2, \\ t - 1 / 4 & \text { for } 1 / 2 \leqslant t \leqslant 3 / 4, \\ 2 t - 1 & \text { for } 3 / 4 \leqslant t \leqslant 1. \end{array} \right.\tag{4}
\]

The function \(\varphi\) is affine on each of the three intervals \([0,1/2]\) , \([1/2,3/4]\) and \([3/4,1]\) ; it is therefore continuous. We have \(\varphi(0) = 0\) and \(\varphi(1) = 1\) . It follows from formula (1) of III, p.~256 defining the juxtaposition of paths that

\[
c _ {1} * (c _ {2} * c _ {3}) = ((c _ {1} * c _ {2}) * c _ {3}) \circ \varphi .
\]

According to the example in III, p.~289, the paths \(c_{1} * (c_{2} * c_{3})\) and \((c_{1} * c_{2}) * c_{3}\) are strictly homotopic, whence equality (3).

Similarly, the function \(\psi\colon\mathbf{I}\to\mathbf{I}\) defined by

\[
\psi (t) = \left\{ \begin{array}{l l} 2 t & \text { for } 0 \leqslant t \leqslant 1 / 2, \\ 1 & \text { for } 1 / 2 \leqslant t \leqslant 1 \end{array} \right.\tag{5}
\]

is continuous and satisfies \(\psi(0)=0,\psi(1)=1\) . The equality \(\gamma_{1}\varepsilon_{y}=\gamma_{1}\) then follows from the fact that

\[
c _ {1} * e _ {y} = c _ {1} \circ \psi
\]

and the equality \(\varepsilon_{x}\gamma_{1}=\gamma_{1}\) is proved similarly, whence (1).

The map \(\sigma\colon\mathbf{I}\times\mathbf{I}\to\mathrm{X}\) defined by

\[
\sigma (t, s) = \left\{ \begin{array}{l l} c _ {1} (2 t s) & \text { for } 0 \leqslant t \leqslant 1 / 2, \\ c _ {1} (2 (1 - t) s) & \text { for } 1 / 2 \leqslant t \leqslant 1 \end{array} \right.\tag{6}
\]

is continuous; it is a strict homotopy connecting the path \(e_{x}\) to the path \(c_{1}*\overline{c_{1}}\) , whence the first equality in (2). The second follows from the first and from the fact that, for every path c, one has \(c=\overline{\overline{c}}\) .

Remark 1. --- Let X be a topological space. Let n be an integer \(\geqslant 1\) and let \((c_{1},\ldots,c_{n})\) be a sequence of paths in X such that \(c_{i}\) and \(c_{i+1}\) are juxtaposable for \(1 \leqslant i \leqslant n-1\) (such a sequence is called a sequence of juxtaposable paths). Denote by c the path

\[
c _ {1} * \left(c _ {2} * \left(\dots * \left(c _ {n - 1} * c _ {n}\right) \ldots\right)\right)
\]

and \(c'\) the path defined by \(c'(t) = c_i(nt - i + 1)\) for \(1 \leqslant i \leqslant n\) and \(t \in [\frac{i-1}{n}, \frac{i}{n}]\) . The paths \(c\) and \(c'\) have the same image and are strictly homotopic: one is the composite of the other with a homeomorphism of I fixing 0 and 1 (cf. III, p. 289, example). We will sometimes denote \(c_{1} * c_{2} * \cdots * c_{n}\) the path \(c'\) .

There exists a unique directed graph \(\varpi(\mathrm{X})\) whose set of vertices is X and whose set of arrows connecting a point \(x\) to a point \(y\) is \(\varpi_{x,y}(\mathrm{X})\) , and in which the maps \(C_{x,y,z}\) define a composition law. By Proposition 1, \(\varpi(\mathrm{X})\) is a groupoid (II, p.~162, def. 4). For every \(x \in \mathrm{X}\) , the identity element at the vertex \(x\) of this groupoid is the class of the constant loop with image \(x\) . The inverse of an arrow \(\gamma\) is the arrow \(\overline{\gamma}\) , which we shall also denote by \(\gamma^{-1}\) . In particular, the composition law \(C_{x,x,x}\) equips, for each \(x \in \mathrm{X}\) , the set \(\varpi_{x,x}(\mathrm{X})\) with a group structure; we denote this group \(\pi_1(\mathrm{X}, x)\) .

DEFINITION 3. --- Let X be a topological space. The groupoid \(\varpi(X)\) is called the Poincaré groupoid, or fundamental groupoid, of the space X. Let x be a point of X; the group \(\pi_1(X,x)\) of loop classes at x is called the Poincaré group, or fundamental group, of the space X at the point x.

Let \(\mathcal{U}\) be a set of subsets of X whose interiors cover X. The path classes in X whose image is contained in one of the subsets belonging to \(\mathcal{U}\) generate the groupoid \(\varpi(\mathrm{X})\) (lemma 4 of III, p.~272).

The orbits of the groupoid \(\varpi(X)\) (II, p.~162) coincide with the arcwise connected components of the space X. In particular (loc. cit.), one has:

PROPOSITION 2. --- Let X be a topological space.

\begin{enumerate}
\def\labelenumi{\alph{enumi})}
\item
  If the space X is path-connected, the groups \(\pi_1(\mathrm{X},x)\) and \(\pi_1(\mathrm{X},y)\) are isomorphic for all points \(x\) and \(y\) of \(\mathrm{X}\) .
\item
  Let x be a point of X; the following conditions are equivalent:
\end{enumerate}

\begin{enumerate}
\def\labelenumi{(\roman{enumi})}
\item
  The group \(\pi_1(X,x)\) is trivial;
\item
  Two paths with origin x in X that have the same endpoint are strictly homotopic;
\item
  Every loop with origin x in X is strictly homotopic to the constant loop with image x.
\end{enumerate}

More precisely, let X be a path-connected topological space and let x and y be points of X. For every element \(\delta\) of \(\varpi_{x,y}(X)\) , the map \(u_{\delta} \colon \gamma \mapsto \delta \gamma \delta^{-1}\) from \(\pi_1(\mathrm{X},y)\) to \(\pi_1(\mathrm{X},x)\) is a group isomorphism whose inverse isomorphism is \(u_{\delta^{-1}}\) . For \(\delta \in \varpi_{x,y}(\mathrm{X})\) and \(\gamma \in \pi_1(\mathrm{X},y)\) we set \(^{\delta}\gamma = u_{\delta}(\gamma)\) . When \(x = y\) , we have \(^{\delta}\gamma = \operatorname{Int}(\delta)(\gamma)\) , i.e.~\(u_{\delta} = \operatorname{Int}(\delta)\) .

Let \(x, y, z\) be points of X. For \(\delta \in \varpi_{x,y}(X)\) and \(\eta \in \varpi_{y,z}(X)\) , we have \(u_{\delta \eta} = u_{\delta} \circ u_{\eta}\) , which is also written \(\delta \eta \gamma = {}^{\delta}(\eta \gamma)\) for \(\gamma \in \pi_1(X, z)\) .

Let \(x, y\) be points of X and let \(\delta\) , \(\delta'\) be elements of \(\varpi_{x,y}(\mathrm{X})\) . We have

\[
u _ {\delta^ {\prime}} = u _ {\delta} \circ \operatorname{Int} (\delta^ {- 1} \delta^ {\prime}) = \operatorname{Int} (\delta^ {\prime} \delta^ {- 1}) \circ u _ {\delta}.\tag{7}
\]

Remarks. --- 2) Let X be a topological space and let C be a path-connected component of X. If x and y are points of C, the topological space \(\Lambda_{x,y}(\mathrm{C})\) is identified with the topological space \(\Lambda_{x,y}(\mathrm{X})\) , so that the set \(\varpi_{x,y}(\mathrm{C})\) is identified with the set \(\varpi_{x,y}(\mathrm{X})\) . Thus, the fundamental group \(\varpi(\mathrm{C})\) is identified with the full subgroupoid of \(\varpi(\mathrm{X})\) having C as its set of points. In particular, for every point x of C, the group \(\pi_1(\mathrm{C}, x)\) is identified with the group \(\pi_1(\mathrm{X}, x)\) .

\begin{enumerate}
\def\labelenumi{\arabic{enumi})}
\setcounter{enumi}{2}
\tightlist
\item
  Let X be a topological space and let \(x\) , \(y\) , \(z\) be points of X. The map \((c,d)\mapsto c*d\) from \(\Lambda_{x,y}(\mathrm{X})\times \Lambda_{y,z}(\mathrm{X})\) to \(\Lambda_{x,z}(\mathrm{X})\) is continuous (III, p.~258, corollary of prop. 2). Note that the composition map \(\varpi_{x,y}(\mathrm{X})\times \varpi_{y,z}(\mathrm{X})\to \varpi_{x,z}(\mathrm{X})\) deduced from it is not necessarily continuous when one equips the sets \(\varpi_{x,y}(\mathrm{X})\) , \(\varpi_{y,z}(\mathrm{X})\) and \(\varpi_{x,z}(\mathrm{X})\) with quotient topologies (cf.~TG, I, p.~35 and III, p.~259). However, for every \(\gamma_0\in \varpi_{x,y}(\mathrm{X})\) and every \(\delta_0\in \varpi_{y,z}(\mathrm{X})\) , the partial mappings \(\gamma \mapsto \gamma \delta_0\) from \(\varpi_{x,y}(\mathrm{X})\) to \(\varpi_{x,z}(\mathrm{X})\) and \(\delta \mapsto \gamma_0\delta\) from \(\varpi_{y,z}(\mathrm{X})\) to \(\varpi_{x,z}(\mathrm{X})\) are homeomorphisms. Indeed, let \(c_0\in \Lambda_{x,y}(\mathrm{X})\) be a path of class \(\gamma_0\) . The map \(d\mapsto c_0*d\) is a continuous map from \(\Lambda_{y,z}(\mathrm{X})\) to \(\Lambda_{x,z}(\mathrm{X})\) . The map \(\delta \mapsto \gamma_0\delta\) follows by passing to quotients, hence is continuous. The same holds for the map \(\delta^{\prime}\mapsto \gamma_{0}^{-1}\delta^{\prime}\) from \(\varpi_{x,z}(\mathrm{X})\) to \(\varpi_{y,z}(\mathrm{X})\) . Since these two maps are inverses of each other, they are homeomorphisms. One argues analogously for the map \(\gamma \mapsto \gamma \delta_0\) . See also IV, p.~374.
\end{enumerate}

